
Semantic Role Labelling~\citep[SRL, ][]{marquez08srl} is generally understood as 
the task of identifying and classifying the semantic arguments and modifiers of 
the predicates mentioned in a sentence. For exampl
e, in the case of the following sentence\footnote{``Haag plays Elianti'' is a 
segment of a sentence in training corpus.}:
\begin{quote}
\begin{center}
    \includegraphics[scale=.63]{haag-example}
\end{center}
\end{quote}
"have to" we are to find out that for the predicate token {}``plays'' with sense 
``play a role'' (play.02) the phrase headed by the token {}``Haag'' is referring 
to the player (A0) of the play event, and the phrase headed by the token 
{}``Elianti''  is referring to the role (A1) being played. SRL is considered as 
a key task for applications that require to answer {}``Who'', {}``What'', 
{}``Where'', etc. questions, such as Information Extraction, Question Answering 
and Summarization. 

In the next sections we will first briefly introduce Markov Logic. Then we 
present the Markov Logic model we used in our final submission. We present and 
analyse our results in section \ref{sec:results} before we conclude in Section 
\ref{sec:conclusion}.

